% Section 1 of the report, standalone so it compiles on its own.
% To move it into the report: copy the \tikzset/\definecolor/\lstset lines into the report preamble once,
% then everything between \begin{document} and the bibliography.
% sec:interface is the label of section 2 (not written yet), so it shows as ?? here.
\documentclass[11pt,a4paper]{article}
\usepackage[utf8]{inputenc}\usepackage[T1]{fontenc}\usepackage{lmodern}
\usepackage[margin=2.2cm]{geometry}
\usepackage{tikz}\usetikzlibrary{positioning,arrows.meta,calc,fit,decorations.pathreplacing,backgrounds}
\usepackage{booktabs,array,amssymb,caption,listings,xcolor,url,amsmath}
\numberwithin{figure}{section}\numberwithin{table}{section}
\AtBeginDocument{\numberwithin{lstlisting}{section}}
\captionsetup{font=small,labelfont=bf}

\definecolor{motion}{HTML}{DCE8F5}
\definecolor{percep}{HTML}{E3F1E0}
\definecolor{outside}{HTML}{F3F3F3}
\definecolor{accent}{HTML}{C0504D}

\tikzset{
  comp/.style={draw, rounded corners=3pt, thick, inner sep=6pt, align=center},
  prog/.style={draw, rounded corners=2pt, fill=white, font=\small, align=center, minimum height=8mm, inner sep=4pt},
  dev/.style={draw, rounded corners=3pt, fill=outside, font=\small, align=center, minimum height=9mm, inner sep=5pt},
  link/.style={-{Stealth[length=2.2mm]}, thick},
  wire/.style={thick},
  lbl/.style={font=\footnotesize, align=center},
}

\lstset{language=Python, basicstyle=\ttfamily\small, columns=fullflexible, keepspaces=true,
  commentstyle=\color{gray}\itshape, keywordstyle=\bfseries, frame=lines, xleftmargin=1em}

\begin{document}

\section{Robot Interface}
\label{sec:commanding}

This section describes what came with the Lite3 robot, and what was missing before the robot could be
controlled from a program. Subsection~\ref{sec:requirements} states the requirements.
Subsection~\ref{sec:access} describes how the robot's computers were reached. Subsections~\ref{sec:bridge}
to~\ref{sec:listening} follow the findings that led to a working way of controlling the robot, and
Subsection~\ref{sec:findings-summary} sums them up.

\subsection{Requirements}
\label{sec:requirements}

The Lite3 was delivered with a control app from the manufacturer, DEEP Robotics. The app runs on a handheld
controller that came with the robot, which is an Android device with two joysticks. The same app can also
run on an ordinary phone. In the rest of this report, both are simply called the app. The robot and the app
are shown in Figure~\ref{fig:robot-app}. With the app, an operator can make the robot stand up and lie
down, walk in any direction, change its gait (the pattern in which the legs move) and speed, and perform
built-in tricks.

\begin{figure}[htbp]
\centering
\fbox{\parbox[c][5cm][c]{0.7\linewidth}{\centering\small Photo: the robot, with the app\\ on the handheld (to be added)}}
\caption{The Lite3 robot and its control app. The supplied handheld is an Android device with joysticks
running the same app as a phone.}
\label{fig:robot-app}
\end{figure}

The goal of this part of the project was to do the same from a program. This gave two requirements. First,
a program must be able to make the robot stand up, lie down and walk, without the app. Second, the speed
and direction of walking must be able to come from an algorithm, such as a navigation system that steers
the robot towards a goal.

No program interface for this came with the robot. Finding out how the robot could be controlled was done
by trial and error, and the work went back and forth between several leads. Each finding narrowed down
what was missing. The findings are presented in the order that best explains how they fit together, which
is close to the order in which they were made.

\subsection{Access to the Robot}
\label{sec:access}

The Lite3 has two computers inside it, as shown in Figure~\ref{fig:system-as-found}. The motion computer
runs the software that makes the legs walk. The perception computer is meant for sensors and for
higher-level programs, such as navigation. Both computers run Linux.

\begin{figure}[htbp]
\centering
\begin{tikzpicture}[node distance=6mm and 10mm]
  % motion computer
  \node[prog] (jy) {Walking controller\\ (\texttt{jy\_exe}, closed)};
  \node[prog, below=6mm of jy] (sdk) {Motion SDK\\ (stand-up demo)};
  % perception computer
  \node[prog, right=40mm of jy] (br) {ROS 2 bridge\\ (velocity in, state out)};
  \node[prog, below=6mm of br] (ros) {ROS 2 programs\\ (later: navigation)};
  \begin{scope}[on background layer]
    \node[comp, fill=motion, fit=(jy)(sdk), inner sep=8pt] (mc) {};
    \node[comp, fill=percep, fit=(br)(ros), inner sep=8pt] (pc) {};
    \node[draw, dashed, rounded corners=6pt, fit=(mc)(pc), inner sep=14pt] (robot) {};
  \end{scope}
  \node[font=\small\bfseries, anchor=south] at (mc.north) {Motion computer};
  \node[font=\small\bfseries, anchor=south] at (pc.north) {Perception computer};
  \node[font=\small, anchor=north] at (robot.south) {Inside the robot};
  % outside the robot
  \node[dev, above=20mm of mc] (app) {The app\\ (handheld or phone)};
  \node[dev, above=20mm of pc] (lap) {Laptop};

  % links inside
  \draw[link] (sdk) -- node[lbl, right] {UDP} (jy);
  \draw[link] ([yshift=2mm]br.west) -- node[lbl, above] {velocity (UDP)} ([yshift=2mm]jy.east);
  \draw[link] ([yshift=-2mm]jy.east) -- node[lbl, below] {state (UDP)} ([yshift=-2mm]br.west);
  \draw[link] (ros) -- (br);
  % app and laptop
  \draw[link] (app) -- node[lbl, left] {Wi-Fi hotspot:\\ commands} (app |- robot.north);
  \coordinate (split) at ($(lap.south)+(0,-8mm)$);
  \draw[wire, dashed] (lap.south) -- (split);
  \draw[link, dashed] (split) -- (split |- robot.north);
  \draw[link, dashed] (split) -| ([xshift=6mm]app.east |- robot.north);
  \node[lbl, anchor=south west] at ([xshift=1mm]split) {SSH, over the hotspot\\ or an Ethernet cable};
\end{tikzpicture}
\caption{The robot as found. The motion computer runs the vendor's closed walking controller and the
vendor's motion SDK; the perception computer runs the vendor's ROS~2 bridge. The app and a laptop reach the
robot over its Wi-Fi hotspot or a cable. Solid arrows: commands and state; dashed: remote access.}
\label{fig:system-as-found}
\end{figure}

Both computers were reached from a laptop using SSH, a standard tool for logging in to another computer
over a network and running commands on it. The laptop was connected either to the robot's own Wi-Fi
hotspot, which is the same network the app uses, or with an Ethernet cable plugged into the robot. With access in place, both computers were searched for
software that could command the robot.

\subsection{The ROS~2 Bridge}
\label{sec:bridge}

On the perception computer, the manufacturer had installed a bridge to ROS~2. ROS~2 (Robot Operating
System~2) is a widely used framework in which programs exchange messages on named channels, called topics.
Many robot tools, including navigation systems, are built on it. A bridge like this was expected to give
full control of the robot from ROS~2, with commands to stand up, lie down and switch between modes.

This turned out not to be the case. The bridge only forwards velocity commands, meaning how fast the robot
should walk forward, sideways and turn, from ROS~2 to the motion computer. It also passes some of the
robot's state back to ROS~2. It has no command for standing up, lying down or anything else.

To test the bridge, the robot was first stood up with the app. A command to walk forward was then sent
through the bridge. The robot started stepping in place, but did not move forward. The command reached the
robot in some form, but something stopped the robot from following it. The reason was found later, and is
explained in Subsection~\ref{sec:auto-mode}. At this point, the bridge alone could not meet either of the
requirements.

\subsection{The Walking Controller and the Motion SDK}
\label{sec:sdk}

On the motion computer, the legs are controlled by the manufacturer's walking controller, a program called
\texttt{jy\_exe}. It runs at all times, receives commands from the app and from the ROS~2 bridge, and turns
them into leg movements. At first, it looked like the place to find the missing commands. However, its source code is
not available, so how it works inside could not be read.

Next to it, the manufacturer had installed a software development kit (SDK) for motion, which is also
published on GitHub~\cite{motionsdk}. The SDK is made for controlling each joint of the legs directly. Making
the robot walk this way would mean writing a complete walking controller from scratch, while the
manufacturer's walking controller already walks well. The aim was therefore to command the existing walking
controller, not to replace it, so the SDK was not a solution in itself. However, its example program
contains a demo that stands the robot up, which is not described in the documentation.

The demo gave an important clue. The SDK talks to the walking controller over the network using UDP, a
simple way of sending short messages between computers. Besides the joint commands, the SDK has a function
for sending a single short command to the walking controller. This fitted with how the app works, since
the app also controls the robot over the network, as shown in Figure~\ref{fig:system-as-found}.

\subsection{Standing Up Without the App}
\label{sec:sdk-stand}

The demo was built and run on the motion computer with the app switched off, and the robot stood up.
Figure~\ref{fig:sdk-demo} shows what the demo sends. It first sends one short command that hands control of
the joints over to the SDK. It then sends joint commands, one message every millisecond, which fold the
legs in, raise the robot and hold it standing. Finally, a second short command hands control back to the
walking controller.

\begin{figure}[htbp]
\centering
\begin{tikzpicture}[x=0.62cm, y=1cm]
  % phases (seconds from the start of the demo)
  \fill[motion] (0,0.3) rectangle (7,1.1);
  \fill[percep!70!motion] (7,0.3) rectangle (8,1.1);
  \fill[percep] (8,0.3) rectangle (9.5,1.1);
  \fill[percep!40] (9.5,0.3) rectangle (17,1.1);
  \draw (0,0.3) rectangle (17,1.1);
  \foreach \x in {7,8,9.5} \draw (\x,0.3) -- (\x,1.1);
  \node[font=\footnotesize, align=center] at (3.5,0.7) {take control,\\ joints to zero};
  \node[font=\footnotesize] at (13.25,0.7) {hold the standing pose};
  \draw (7.5,1.1) -- (7.5,1.55) node[above, font=\footnotesize] {legs in};
  \draw (8.75,1.1) -- (8.75,2.15) node[above, font=\footnotesize] {stand up};
  % command packets
  \draw[accent, line width=1.6pt] (0,1.1) -- (0,1.9) node[above, font=\footnotesize, align=center, text=black] {command:\\ take control};
  \draw[accent, line width=1.6pt] (17,1.1) -- (17,1.9) node[above, font=\footnotesize, align=center, text=black] {command:\\ hand control back};
  % time axis
  \draw[->, thick] (0,0) -- (18.2,0) node[right, font=\small] {$t$ [s]};
  \foreach \t in {0,5,10,15} \draw (\t,0.1) -- (\t,-0.1) node[below, font=\footnotesize] {\t};
  \draw[decorate, decoration={brace, mirror, amplitude=5pt}] (7,-0.55) -- (17,-0.55)
     node[midway, below=6pt, font=\footnotesize] {joint commands, one message every millisecond};
  % afterwards
  \node[draw, rounded corners=2pt, fill=outside, font=\footnotesize, align=left, anchor=west] at (19.9,0.9)
     {Afterwards:\\ still standing,\\ walks on the app's\\ joysticks, ignores\\ the ROS~2 bridge};
\end{tikzpicture}
\caption{What the vendor's stand-up demo sends. One short command hands the joints to the SDK, joint
commands then raise the legs and hold them, and a second short command hands control back to the walking
controller. Durations are approximate, from the demo's code.}
\label{fig:sdk-demo}
\end{figure}

This was half a success. The robot stood up without the app, and it kept standing after the demo had
ended. When the app was connected afterwards, the robot walked on the app's joysticks, but it still ignored
velocity commands from the ROS~2 bridge.

The demo also showed that the walking controller accepts short commands that are not described anywhere.
Handing control over and taking it back were each done with a single message. This suggested that there
were more commands of this kind, and that the app was most likely sending them over the network.

\subsection{Auto Mode}
\label{sec:auto-mode}

While trying out the app, a menu with a setting called auto mode was found by chance. With auto mode
switched on, the same forward command as in Subsection~\ref{sec:bridge} made the robot walk forward. At
first, it was not clear why.

The explanation was that the robot had been in manual mode. In manual mode, the robot follows only the
app's joysticks and ignores velocity commands from anywhere else. This is why the robot stepped in place in
Subsection~\ref{sec:bridge}, and why it ignored the bridge after the demo in
Subsection~\ref{sec:sdk-stand}. The robot does not report which mode it is in, so a program cannot check
it. A program therefore has to switch on auto mode itself before it sends any velocity commands.

With this, the robot could walk on commands from a program. However, it only worked while the app was
connected to the robot.

\subsection{The Missing Piece}
\label{sec:missing}

The robot could now be stood up from a program, and velocity commands worked in auto mode. Both still
needed the app to be connected. When the app was switched off, the robot stopped accepting velocity
commands, so the app had to be providing something that was still missing.

The documentation of the motion SDK gave a hint. It states that if the walking controller receives no
commands from the SDK for more than one second, it takes back control of the robot~\cite{motionsdk}. This
suggested that the walking controller might expect to hear from whatever controls it all the time, and not
only when something changes. If the app sent such a continuous signal, it would be visible on the
network.

\subsection{Listening to the Network}
\label{sec:listening}

To find this signal, the network traffic on the motion computer was recorded with tcpdump, a standard tool
that records the messages passing through a computer's network connections~\cite{tcpdump}. The
documentation of the motion SDK uses tcpdump in the same way for troubleshooting. While recording, buttons
such as stand were pressed in the app.

Figure~\ref{fig:capture} shows such a recording, in which auto mode was switched on and stand was pressed
twice. Three kinds of traffic can be seen in the figure.
First, the app sends the same short message every half second, also when nothing is pressed. This is a
heartbeat, a message whose only purpose is to tell the walking controller that the app is still
connected. Second, each button press shows up as one single message. Switching on auto mode is a message of its
own, which explains the discovery in Subsection~\ref{sec:auto-mode}. Pressing stand sent the same message
both times; the robot stood up on the first and lay down on the second, so the stand command switches
between standing and lying. Third, the robot sends its state back. The app receives a few status messages
per second, while the perception computer receives a steady stream of about 2000 messages per second, which
the ROS~2 bridge reads.

\begin{figure}[htbp]
\centering
\begin{tikzpicture}[x=0.175cm, y=1cm]
  % real capture on the motion computer, 2026-11-23; times in seconds from the start of the recording
  \foreach \y/\name in {3/{App $\rightarrow$ robot}, 1.9/{Robot $\rightarrow$ app}, 0.9/{Robot $\rightarrow$ perception\\ computer}} {
    \draw[gray] (0,\y) -- (65,\y);
    \node[anchor=east, font=\small, align=right] at (-0.8,\y) {\name};
  }
  % the app's heartbeat
  \foreach \t in {0.01,0.52,1.01,1.51,2.01,2.51,3.01,3.51,4.01,4.52,5.01,5.51,6.01,6.51,7.01,7.51,8.01,8.51,9.01,9.51,10.01,10.51,11.01,11.51,12.01,12.51,13.01,13.51,14.01,14.51,15.01,15.51,16.01,16.52,17.01,17.51,18.01,18.51,19.01,19.51,20.01,20.51,21.01,21.51,22.01,22.51,23.01,23.51,24.01,24.51,25.01,25.51,26.01,26.51,27.01,27.51,28.01,28.51,29.02,29.51,30.01,30.51,31.01,31.51,32.01,32.51,33.01,33.51,34.01,34.51,35.01,35.51,36.01,36.51,37.01,37.51,38.01,38.51,39.01,39.51,40.01,40.51,41.01,41.51,42.01,42.51,43.01,43.51,44.01,44.52,45.01,45.51,46.01,46.51,47.01,47.51,48.01,48.52,49.01,49.51,50.02,50.52,51.02,51.51,52.01,52.51,53.01,53.51,54.01,54.51,55.01,55.52,56.01,56.51,57.01,57.51,58.01,58.51,59.01,59.51,60.01,60.52,61.01,61.51,62.01,62.52,63.01,63.52,64.01,64.51} \draw (\t,3) -- (\t,3.3);
  % the three button presses, one packet each
  \draw[accent, line width=1.4pt] (13.39,3) -- (13.39,3.85) node[above, font=\footnotesize, text=black] {auto mode on};
  \draw[accent, line width=1.4pt] (28.30,3) -- (28.30,3.85) node[above, font=\footnotesize, text=black] {stand: stands up};
  \draw[accent, line width=1.4pt] (42.76,3) -- (42.76,3.85) node[above, font=\footnotesize, text=black] {stand: lies down};
  % the robot's status messages to the app
  \foreach \t in {0.20,0.70,0.70,0.70,0.70,0.70,1.20,1.50,1.50,1.70,2.20,2.30,2.30,2.70,3.10,3.10,3.20,3.20,3.70,3.90,3.90,4.20,4.70,4.70,4.70,5.20,5.50,5.50,5.70,5.70,5.71,6.20,6.30,6.30,6.70,7.10,7.10,7.20,7.70,7.90,7.90,8.20,8.20,8.70,8.70,8.70,9.20,9.50,9.50,9.70,10.20,10.30,10.30,10.70,10.70,10.70,11.10,11.10,11.20,11.70,11.90,11.90,12.20,12.70,12.70,12.70,13.20,13.20,13.50,13.50,13.70,14.20,14.30,14.30,14.70,15.10,15.10,15.20,15.70,15.70,15.71,15.90,15.90,16.20,16.70,16.70,16.70,17.20,17.50,17.50,17.70,18.20,18.20,18.30,18.30,18.70,19.10,19.10,19.20,19.70,19.90,19.90,20.20,20.70,20.70,20.70,20.70,20.70,21.20,21.50,21.50,21.70,22.20,22.30,22.30,22.70,23.10,23.10,23.20,23.20,23.70,23.90,23.90,24.20,24.70,24.70,24.70,25.20,25.50,25.50,25.70,25.70,25.70,26.20,26.30,26.30,26.70,27.10,27.10,27.20,27.70,27.90,27.90,28.20,28.20,28.70,28.70,28.70,29.20,29.50,29.50,29.70,30.20,30.30,30.30,30.70,30.70,30.70,31.10,31.10,31.20,31.70,31.90,31.90,32.20,32.70,32.70,32.70,33.20,33.20,33.50,33.50,33.70,34.20,34.30,34.30,34.70,35.10,35.10,35.20,35.70,35.70,35.70,35.90,35.90,36.20,36.70,36.70,36.70,37.20,37.50,37.50,37.70,38.20,38.20,38.30,38.30,38.70,39.10,39.10,39.20,39.70,39.90,39.90,40.20,40.70,40.70,40.70,40.70,40.70,41.20,41.50,41.50,41.70,42.20,42.30,42.30,42.70,43.10,43.10,43.20,43.20,43.70,43.90,43.90,44.20,44.70,44.70,44.70,45.20,45.50,45.50,45.70,45.70,45.70,46.20,46.30,46.30,46.70,47.10,47.10,47.20,47.70,47.90,47.90,48.20,48.20,48.70,48.70,48.70,49.20,49.50,49.50,49.70,50.20,50.30,50.30,50.70,50.70,50.70,51.10,51.10,51.20,51.70,51.90,51.90,52.20,52.70,52.70,52.70,53.20,53.20,53.50,53.50,53.70,54.20,54.30,54.30,54.70,55.10,55.10,55.20,55.70,55.70,55.71,55.90,55.90,56.20,56.70,56.70,56.70,57.20,57.50,57.50,57.70,58.20,58.20,58.30,58.30,58.70,59.10,59.10,59.20,59.70,59.90,59.90,60.20,60.70,60.70,60.70,60.70,60.71,61.20,61.50,61.50,61.70,62.20,62.30,62.30,62.70,63.10,63.10,63.20,63.20,63.70,63.90,63.90,64.20,64.70,64.70,64.70} \draw (\t,1.9) -- (\t,2.15);
  % the state stream to the perception computer, far too dense to draw packet by packet
  \fill[gray!45] (0,0.9) rectangle (65,1.15);
  % labels
  \node[font=\footnotesize, anchor=north west] at (0,2.9) {heartbeat, every 0.5\,s};
  \node[font=\footnotesize, anchor=north west] at (0,1.8) {status, a few per second};
  \node[font=\footnotesize, anchor=north west] at (0,0.8) {state, about 2000 messages a second};
  % time axis
  \draw[->, thick] (0,0) -- (67,0) node[right, font=\small] {$t$ [s]};
  \foreach \t in {0,10,...,60} \draw (\t,0.1) -- (\t,-0.1) node[below, font=\footnotesize] {\t};
\end{tikzpicture}
\caption{A recording of the network traffic on the motion computer while the app was used: auto mode switched on, then
stand pressed twice. Each press is a single message; between presses the app repeats a heartbeat every half
second. The robot answers the app with a few status messages per second, and streams its full state to the
perception computer.}
\label{fig:capture}
\end{figure}

The last step was to test this with the app switched off. A program sent the heartbeat itself, followed by
the stand command, and the robot stood up. The program then switched on auto mode and sent a forward
velocity command through the ROS~2 bridge, and the robot walked forward. The heartbeat was the missing
piece.

Together, these findings give a complete recipe, shown as simplified code in Listing~\ref{lst:recipe}. A
program keeps sending the heartbeat in the background, sends the stand command and waits until the robot
reports that it is standing, switches on auto mode, and then sends velocity commands through the ROS~2
bridge. The names in the listing are placeholders, but the steps and their order are the same as in the
interface described in Section~\ref{sec:interface}.

\begin{lstlisting}[float=htbp, captionpos=b, label=lst:recipe,
  caption={The recipe found in this section, as simplified code. The names are placeholders; the steps and their order are those of the interface.}]
start_background(send_every(0.5, HEARTBEAT))   # keeps the controller listening

send(STAND)                                    # a toggle, not "stand up"
wait_until(lambda: robot_state().posture == STANDING)

send(AUTO_MODE)                                # otherwise velocity is ignored
while driving:
    publish_velocity(forward=0.3, turn=0.0)    # through the ROS 2 bridge
\end{lstlisting}

\subsection{Summary of the Findings}
\label{sec:findings-summary}

Table~\ref{tab:findings} sums up what each finding added, and what was still missing afterwards. Together,
the findings meet both requirements: a program can stand the robot up, switch it to auto mode and steer it,
without the app.

\begin{table}[htbp]
\centering
\caption{What each finding added, and what was still missing.}
\label{tab:findings}
\begin{tabular}{@{}lll@{}}
\toprule
Finding & Added & Still missing \\
\midrule
ROS~2 bridge   & Velocity reaches the robot & It is ignored \\
Motion SDK     & Stand up without the app   & Velocity still ignored \\
Auto mode      & Walk on velocity           & Needs the app running \\
Heartbeat      & No app needed             & Nothing \\
\bottomrule
\end{tabular}
\end{table}

The recipe in Listing~\ref{lst:recipe} is still fragile. Since the stand command switches between standing
and lying, sending it at the wrong time does the opposite of what was intended. The first stand command
after the robot is switched on is ignored. If the program crashes while the robot is walking,
nothing in the recipe stops it, and the robot only stops by itself after about one and a half seconds. Section~\ref{sec:interface} describes the interface that was built on top of these
findings to handle problems like these.

\begin{thebibliography}{9}
\bibitem{motionsdk} DEEP Robotics, \emph{Lite3\_MotionSDK}, GitHub repository.
  \url{https://github.com/DeepRoboticsLab/Lite3_MotionSDK}
\bibitem{tcpdump} The Tcpdump Group, \emph{tcpdump and libpcap}. \url{https://www.tcpdump.org}
\end{thebibliography}

\end{document}
